% !TEX program = lualatex
\documentclass[aspectratio=43,professionalfonts,handout]{beamer}
\usepackage{beamer_template}

\newcommand{\LectureNum}{5}
\newcommand{\LectureTitle}{微分方程式への応用}
\newcommand{\TextPages}{p.\,62-65}
\newcommand{\CourseName}{応用数学}
\renewcommand{\TermName}{前期}

\begin{document}

\begin{frame}{\CourseName\quad 第\LectureNum 回「\LectureTitle」}
  {\small \TermName\quad \TeacherName\quad ／\quad 教科書 \TextPages}

  \textbf{目標}：ラプラス変換を用いて線形微分方程式の初期値問題・境界値問題を解くことができる

\end{frame}

\begin{frame}{解法の流れ}
  \begin{block}{}
  \begin{itemize}
    \item 微分方程式 $ \to $ ラプラス変換 $ \to $ 代数方程式
    \item 代数方程式を解く $ \to $ 逆ラプラス変換 $ \to $ 解
    \item 初期条件を自動的に組み込める利点
  \end{itemize}
  \end{block}
\end{frame}

\begin{frame}{初期値問題（1階）}
  \begin{block}{}
  \begin{itemize}
    \item 例題1：$dx/dt + x = e^t, x(0)=1$
    \item $\mathcal{L}[x'] = sX(s) - x(0)$ の活用
    \item $X(s)$ の計算 $ \to $ 逆変換
  \end{itemize}
  \end{block}
\end{frame}

\begin{frame}{初期値問題（2階）}
  \begin{block}{}
  \begin{itemize}
    \item 例題2：$x'' + 4x = e^{-t}$
    \item $\mathcal{L}[x''] = s^{2}X(s) - sx(0) - x'(0)$
    \item 部分分数分解 $ \to $ 逆変換
  \end{itemize}
  \end{block}
\end{frame}

\begin{frame}{境界値問題・一般解}
  \begin{block}{}
  \begin{itemize}
    \item 例題3：境界条件の扱い
    \item 例題4：一般解（任意定数を含む解）
  \end{itemize}
  \end{block}
\end{frame}

\begin{frame}{演習}
  \begin{exampleblock}{自分で解いてみよう}
  \begin{itemize}
    \item 問1〜4より選題
  \end{itemize}
  \end{exampleblock}
\end{frame}

\begin{frame}{まとめ}
  \begin{itemize}
    \item ラプラス変換による微分方程式の解法手順
    \item 次回：たたみ込み
  \end{itemize}
\end{frame}

\end{document}
